\section{Data}
\label{sec:data}

Six public sources, all fetched by script rather than vendored, and named in
full because the benchmark's credibility rests on them.

\paragraph{Building Data Genome Project 2} \citep{miller2020bdg2}. $1{,}636$
metered non-residential buildings across nineteen sites, of which $1{,}578$
carry an electricity meter, hourly with matched weather, two years. It is the
only harmonised multi-country building corpus in existence, which is why it
carries the controlled comparison: the same physical quantity, the same metering
convention, the same period. It contains \emph{no Indian buildings}, which is
why the next source is in the study.

\paragraph{I-BLEND} \citep{rashid2019iblend}. One-minute mains power and power
factor for seven buildings on the IIIT-Delhi campus --- academic, lecture,
library, facilities, dining and two dormitories --- and for the three
transformers that feed the whole campus, from August 2013 to December 2017,
CC0. The billing blocks are formed by averaging sixty one-minute readings, not
by interpolating up from hourly data, so this is the only building-level
source in the study at the cadence Indian demand charges are assessed on. Its
documentation states, and the load--temperature signature of each meter
confirms, that every building except Facilities is served by a central chiller
plant that is \emph{not} on the building's own meter: at building level the
corpus is therefore a forecasting and calibration benchmark rather than a
control one, and the selection rule's chilled-water clause records it as such.
The campus feed is the exception, and it is the more important object: an
Indian campus is billed as a single high-tension consumer with one contract
demand and one monthly demand charge, so the transformer total --- median
$199$\,kW, peak $842$\,kW, chiller included --- is what the tariff actually
sees. Meter uptime varies by feed; each series' admissible test Junes and the
training months behind each are recorded by \texttt{data/iblend.py} and
printed in the table.

\paragraph{Delhi State Load Despatch Centre.} Fifteen-minute metered demand for
the city of Delhi, disaggregated across five distribution companies (BRPL, BYPL,
NDPL, NDMC and MES) plus the total, two complete years at full coverage. With
I-BLEND this gives three levels of aggregation --- building, campus, city --- in
one city, at native resolution, under one weather feed. Every non-Indian series
in the study is interpolated up from hourly.

\paragraph{ENTSO-E transmission data, via Open Power System Data.} Hourly
national load for the United Kingdom, Ireland, Germany, France and Spain, drawn
from the European transmission system operators' transparency platform.

\paragraph{Chinese provincial demand} \citep{wu2023china}. Hourly load for the
31 provinces of mainland China, 2018, from which six are drawn. This source was
added after the study was otherwise complete, and it is the only tier-2 series
that is constructed rather than metered. Section~\ref{sec:china} audits it.

\paragraph{ERA5 reanalysis, via the Open-Meteo archive.} Hourly temperature, dew
point and cloud cover at each location, used identically for every series.

\paragraph{Windows.} The building tier and the European national tier share
2016--2017. The Delhi archive covers 2011--2012, so the same protocol is shifted
back five years with its shape preserved. The Chinese series is a single
calendar year, 2018. In every row of the study the test month is June, and the
headline metric is skill against each series' own seasonal-naive baseline: a
ratio computed inside one series' own data and its own period, which is what
makes rows with different units and different years comparable at all. Absolute
losses are not comparable across groups and are never compared.

\paragraph{Selection.} One rule, stated in advance and applied uniformly: no
district chilled-water meter, $>$$97\%$ coverage, median load 40--800\,kW,
positive load--temperature correlation; ties within a (site, usage) cell broken
by highest $p99/\text{median}$. Thirty supplies pass. Two facts about the rule
were discovered by applying it rather than anticipated when writing it, and both
are reported in Section~\ref{sec:calibration}.
